\documentclass[11pt,a4paper]{article}

\usepackage[margin=25mm]{geometry}
\usepackage[T1]{fontenc}
\usepackage{lmodern}
\usepackage{microtype}
\usepackage{booktabs}
\usepackage{longtable}
\usepackage{enumitem}
\usepackage{xcolor}
\usepackage{hyperref}
\usepackage{listings}
\usepackage{fancyhdr}

\hypersetup{colorlinks=true,linkcolor=blue!55!black,urlcolor=blue!55!black}
\setlist[itemize]{noitemsep,topsep=3pt}
\setlist[enumerate]{topsep=3pt}
\pagestyle{fancy}
\fancyhf{}
\lhead{ARUN-X}
\rhead{Preprocessing and readiness report}
\cfoot{\thepage}

\title{\textbf{ARUN-X Solar Flare Dataset}\\[3pt]
\large Preprocessing from Data Acquisition to Model Readiness}
\author{Project pipeline documentation}
\date{24 November 2026}

\begin{document}
\maketitle

\begin{abstract}
This report describes the ARUN-X data workflow for preparing HEL1OS hard X-ray and SoLEXS soft X-ray measurements with GOES flare-event labels. It follows the data from acquisition and date selection through time alignment, quality handling, label creation, and the planned training/evaluation stage. The repository currently contains a one-day processing example (6 October 2026), coverage-analysis utilities and reports, and label-generation scripts. It does not yet contain a completed multi-day training corpus, finalized chronological splits, or measured model results. Implemented behavior and planned work are distinguished throughout.
\end{abstract}

\tableofcontents
\newpage

\section{Project objective and data sources}

ARUN-X aims to build a reliable time-series dataset for detecting active M/X-class solar flares and forecasting their starts. Three sources are involved:
\begin{itemize}
  \item \textbf{HEL1OS}: hard X-ray light-curve channels in FITS files, with detector good-time intervals (GTIs).
  \item \textbf{SoLEXS}: soft X-ray SDD2 light-curve data and a corresponding GTI.
  \item \textbf{GOES}: an event catalogue containing flare start, peak, end, and class information. It provides labels rather than instrument input features.
\end{itemize}

The intended processing flow is:
\begin{center}
Raw instrument files $+$ GOES event catalogue $\longrightarrow$ coverage/date selection $\longrightarrow$ UTC quality-filtered minute table $\longrightarrow$ labels $\longrightarrow$ chronological splits $\longrightarrow$ model training and evaluation.
\end{center}

\section{Acquisition and date selection}

\subsection{What is present}
The repository contains downloaded HEL1OS and SoLEXS source data, a GOES event-list file, and scripts/reports for auditing observation dates and overlap. Coverage tools inspect HEL1OS segments/GTIs, SoLEXS dates, shared coverage, and GOES event dates. The final-date report combines instrument coverage and event information. The project plan records a proposed February--July 2026 period and illustrative month-based splits, subject to confirming availability from both instruments.

\subsection{What is missing at acquisition}
The listed repository files do not include a maintained, end-to-end bulk-download script for both missions. The one-day source data were obtained, but bulk acquisition is marked as not started in \texttt{z\_plan.md}. A reproducible acquisition step should record the portal/query, product version, requested dates, successful/failed downloads, checksums or file sizes, and extraction status. Date selection must be based on actual common coverage and event counts, rather than the proposed range alone.

\section{GOES event parsing}

The script \texttt{scripts/pipeline/parse\_goes.py} reads the GOES event-list text using a regular expression and writes \texttt{output/goes\_events\_parsed.csv}. It parses date, start, peak, end, and class (for example, M or X plus intensity), converts event times to UTC, and derives \texttt{class\_type} and numeric \texttt{intensity} fields. This table is the source of active-flare and forecast targets.

The parser currently assumes a specific text layout and does not itself document or validate catalogue version/provenance, malformed rows, duplicate events, or time-zone conventions. These checks should be added before scaling to the final dataset.

\section{Instrument preprocessing and merging}

The current one-day implementation is in \texttt{scripts/pipeline/merge\_lightcurves.py}. Its configured demonstration date is 2026-10-06.

\subsection{HEL1OS}
The script finds date-matching \texttt{lightcurve\_*.fits} files. For each detector it locates the corresponding GTI FITS file, reads MJD timestamps and \texttt{CTR} count rates from the light-curve HDUs, and converts MJD to timezone-aware UTC timestamps. Measurements are retained only when finite and within a GTI; otherwise the signal is set to NaN. Detector/channel series are aligned by timestamp. Exact duplicate timestamps within a channel trigger an error.

\subsection{SoLEXS}
The script reads the SDD2 \texttt{RATE} extension, using \texttt{TIME} as Unix seconds and \texttt{COUNTS} as the signal, then converts timestamps to UTC. It reads the matching GTI extension and retains finite counts only inside valid intervals. The present implementation expects exactly one light-curve file and one GTI file under the source tree, which is a one-day assumption that must be replaced for multi-day runs.

\subsection{Resampling and alignment}
The two instrument frames are outer-joined on timestamp. For each numeric signal, valid samples are aggregated as a one-minute mean. A minute's GTI flag is true if at least one sample in that minute falls in a valid GTI. The result is reindexed to a complete UTC day with 1,440 minute timestamps. Signal values missing after aggregation remain NaN; missing flags are true and absent GTI flags false. Values outside GTIs are explicitly set to NaN.

The merged table is written as \texttt{output/merged\_20261006.parquet}. The script checks row count, unique timestamps, unique columns, the expected 21 signal columns, presence of missing/GTI flags, and consistency between flags and values. These validations apply to the configured one-day run; multi-day boundary behavior is not yet established.

\section{Label generation}

\subsection{Active event, class, and peak labels}
\texttt{scripts/pipeline/label\_flares.py} joins the merged observations to parsed GOES M/X events. It creates:
\begin{itemize}
  \item \texttt{flare\_active}: 1 for timestamps from event start through event end, inclusive; otherwise 0.
  \item \texttt{flare\_class}: 0 for no active M/X event, 1 for M, and 2 for X. For overlapping event intervals, the larger class value is retained.
  \item \texttt{flare\_peak}: 1 on the minute obtained by flooring the event peak timestamp to a minute; otherwise 0.
\end{itemize}
The output is \texttt{output/labeled\_20261006.parquet}.

\subsection{Forecast targets}
\texttt{scripts/pipeline/create\_forecast\_labels.py} creates \texttt{flare\_15m}, \texttt{flare\_30m}, and \texttt{flare\_60m}. A row is positive if an M/X event starts strictly after its timestamp and within the selected future horizon. Rows during an active M/X event are assigned $-1$, as are rows too close to the end of the data to have a complete future window. Other rows are assigned 0. The output is \texttt{output/forecast\_labeled\_20261006.parquet}.

The $-1$ value is an exclusion/unknown sentinel and must not be treated as a negative label during training or scoring. Boundary behavior, overlapping events, rows with missing instrument coverage, and catalogue/event time conventions require an explicit multi-day label audit.

\section{From labeled table to model-ready dataset}

The current Parquet files are intermediate labeled data; a complete model-ready dataset is not yet present. The planned last preprocessing stages are:
\begin{enumerate}
  \item Process all selected dates with reproducible per-date inputs and outputs, then concatenate without duplicate or missing UTC minutes at boundaries.
  \item Apply documented quality rules, retaining GTI and missingness indicators and identifying dates/intervals with insufficient instrument coverage. Do not silently replace missing measurements with zero.
  \item Generate and validate all task labels and agreed forecast horizons across the entire date range. SEP-contamination flags should only be included when a reliable source is available.
  \item Construct features from the prediction timestamp and its past only. Candidate features include current signal, changes, and trailing summaries over fixed windows. Exclude GOES-derived labels, future measurements, and target columns from model inputs.
  \item Split by chronology into train, validation, and test dates. Fit imputation, transforms, scaling, feature selection, and any sampling on the training set only; reuse fitted preprocessing on later splits.
  \item Persist a schema and manifest describing dates, source products, quality decisions, target definitions, split boundaries, and preprocessing configuration.
\end{enumerate}

\section{Training and evaluation: planned, not completed}

The recommended first serious model is a class-weighted gradient-boosted tree (CatBoost or LightGBM) on past-only engineered features, compared with a climatology/no-skill baseline and logistic regression. A temporal CNN or CNN-LSTM is a later experiment once substantially more independent events have been collected. This is a recommendation, not a result: the repository has no finalized multi-day splits or trained-model scores.

Evaluation should use untouched later dates, preserve natural class prevalence in validation/test, and report event counts. Accuracy alone is not useful for a rare-event problem. Report precision-recall AUC, threshold-specific precision and recall, event-level recall, false alarms per day, and lead time. Threshold selection and any hyperparameter tuning must use validation data, not the test set. Any X-class-specific claim must account for the small number of X events noted in the project plan.

\section{Missing work and readiness checklist}

\begin{longtable}{p{0.26\textwidth}p{0.22\textwidth}p{0.44\textwidth}}
\toprule
\textbf{Area} & \textbf{Status} & \textbf{Needed before final training} \\
\midrule
Bulk acquisition & Not complete & Reproducible HEL1OS/SoLEXS downloader or logged acquisition process, provenance, integrity checks, and a verified list of successful dates. \\
Date selection & In progress & Confirm common instrument coverage for the chosen period and recount events on eligible dates; finalize split boundaries based on actual data. \\
Multi-day merge & Missing & Generalize the configured one-day date and single-SoLEXS-file assumptions; handle daily boundaries and reruns deterministically. \\
Quality policy & Partial & Define minimum coverage, missing-data handling, bad intervals, instrument anomalies, and any reliable SEP-contamination flags. \\
Label audit & Partial & Validate time conventions, inclusive interval semantics, peak rounding, overlaps, event-to-row matching, and the $-1$ exclusion behavior at scale. \\
Feature pipeline & Missing & Implement reproducible, past-only features and training-fitted transforms; persist schema and transform state. \\
Chronological splits & Missing & Create train/validation/test datasets by date and report valid coverage and independent event counts for each. \\
Model training & Not started & Train baselines and candidate models using training data only; tune on validation and retain the test set for final evaluation. \\
Evaluation & Not started & Produce measured rare-event and event-level metrics with thresholds, lead times, false alarms, and date-level context. \\
Application integration & Not started & Load actual data/model, persist preprocessing with the model, and connect predictions to the application only after pipeline validation. \\
\bottomrule
\end{longtable}

\section{Conclusion}

ARUN-X already demonstrates the central one-day transformation: parse GOES events, read HEL1OS and SoLEXS light curves, enforce GTI and finite-value constraints, align observations on a one-minute UTC grid, and generate active/class/peak plus forecast labels. The next critical deliverable is a verified multi-day dataset with documented quality rules and chronological splits. Until that exists and is evaluated on held-out dates, model selection remains a recommendation rather than an empirical conclusion.

\end{document}
